\chapter{Селектор кубического incidence-назначения двух источников}
\label{ch:v8-extra-mass-cubic-incidence-selector}

\section{Шесть полноранговых кандидатов}

Предыдущий гейт оставил упорядоченные назначения
\begin{equation}
 (e_B,e_M)\in
 \{(Y,u),(Y,d),(u,Y),(u,d),(d,Y),(d,u)\},
 \label{eq:v8-cubic-incidence-selector-menu}
\end{equation}
где $T_B=3$, $T_M=2$ и
\begin{equation}
 q_Y=(0,-3),\qquad q_u=(1,1),\qquad q_d=(-1,1).
 \label{eq:v8-cubic-incidence-selector-charges}
\end{equation}
Для сравнения формы используем inherited trace-метрику
\begin{equation}
 K=\operatorname{diag}(12,48),
 \qquad
 C(e_B,e_M)=\begin{pmatrix}3q_{e_B}&2q_{e_M}\end{pmatrix}.
 \label{eq:v8-cubic-incidence-selector-weighted-map}
\end{equation}

\section{Determinant и метрическая площадь}

Максимизация невзвешенного determinant даёт
\begin{equation}
 |\det Q|_{\max}=3,
 \qquad
 \mathcal M_{\det}=\{(Y,u),(Y,d),(u,Y),(d,Y)\}.
 \label{eq:v8-cubic-incidence-selector-det-survivors}
\end{equation}
Тот же набор получается из площади в двойственной trace-метрике:
\begin{equation}
 \mathcal A(e_B,e_M)=
 \sqrt{\det\bigl(C^TK^{-1}C\bigr)},
 \qquad
 \mathcal A_{Y\text{-pair}}=\frac34,
 \quad \mathcal A_{ud\text{-pair}}=\frac12.
 \label{eq:v8-cubic-incidence-selector-metric-area}
\end{equation}
Следовательно, determinant и объём не дают единственного назначения.

Знак determinant также не является селектором. При замене ориентированного
базиса $A\mapsto-A$ он меняется:
\begin{equation}
 \det Q\longmapsto-\det Q,
 \qquad |\det Q|\longmapsto|\det Q|.
 \label{eq:v8-cubic-incidence-selector-orientation-sign}
\end{equation}

\section{Обусловленность и совпадение кратностей}

Число обусловленности измерим после whitening матрицей $K^{-1/2}$:
\begin{equation}
 \kappa(e_B,e_M)=
 \sqrt{\frac{\lambda_{\max}(C^TK^{-1}C)}
 {\lambda_{\min}(C^TK^{-1}C)}}.
 \label{eq:v8-cubic-incidence-selector-condition-number}
\end{equation}
Минимум достигается ровно на двух назначениях:
\begin{equation}
 \mathcal M_{\kappa}=\{(u,Y),(d,Y)\},
 \qquad
 \kappa_{\min}=
 \sqrt{\frac{27+3\sqrt{17}}{27-3\sqrt{17}}}.
 \label{eq:v8-cubic-incidence-selector-condition-minima}
\end{equation}
Это же меню поддерживает простое совпадение комплексных кратностей:
$T_B=3$ с цветным сектором размерности $3$, а $T_M=2$ со слабым сектором
размерности $2$:
\begin{equation}
 \mathcal M_{\dim}=\{(u,Y),(d,Y)\}.
 \label{eq:v8-cubic-incidence-selector-dimension-match}
\end{equation}
Таким образом, два независимых внутренних критерия согласованы, но оставляют
бинарную свободу $u/d$.

\section{Почему бинарная свобода не снимается}

Для одного общего кубического коэффициента два финалиста дают source-лучи
\begin{equation}
 (u,Y):\ (j_A,j_B)\parallel(3,-3),
 \qquad
 (d,Y):\ (j_A,j_B)\parallel(-3,-3).
 \label{eq:v8-cubic-incidence-selector-final-rays}
\end{equation}
Их щели соответственно лежат на координатных осях:
\begin{equation}
 (\Delta_u,\Delta_d)\parallel(0,-1),
 \qquad
 (\Delta_u,\Delta_d)\parallel(-1,0).
 \label{eq:v8-cubic-incidence-selector-gap-rays}
\end{equation}
Поэтому выбор одного кандидата буквально выбирает, какой цветной канал
получает щель.

Сектора $u$ и $d$ имеют различные гиперзаряды, так что они не связаны
сохранённым gauge-интертвейнером. Но отсутствие симметрии между объектами не
есть правило выбора одного из них. Классификация конечной спектральной
тройки требует явных матриц и диаграммы связей
\cite{v8-cubic-incidence-selector-krajewski,v8-cubic-incidence-selector-cacic};
она не разрешает добавлять ориентированное ребро по меньшему condition
number после просмотра результата.

Equal-gap требование $j_A=0$ не выбирает ни одного финалиста при одном
коэффициенте. При двух коэффициентах оно уничтожает $B$-компоненту:
\begin{equation}
 j_A=\mp9c_B=0\quad\Longrightarrow\quad c_B=0.
 \label{eq:v8-cubic-incidence-selector-equal-gap-ratio}
\end{equation}
Обе ветви тогда схлопываются к одному rank-one назначению $M\mapsto Y$.
Наблюдаемая мишень поэтому не является внутренним incidence-селектором.

Итоговые реестры:
\begin{equation}
 N_{\rm exact\ classification}=8/8,
 \qquad N_{\rm intrinsic\ selector}=0/8,
 \qquad |\mathcal M_{\rm final}|=2.
 \label{eq:v8-cubic-incidence-selector-ledgers}
\end{equation}

\begin{theorem}[Двухточечный остаток incidence-меню]
Максимальный determinant и максимальная trace-метрическая площадь оставляют
четыре назначения. Минимизация точного числа обусловленности и совпадение
кратностей независимо сокращают их до $(u,Y)$ и $(d,Y)$. Все проверенные
внутренние критерии инвариантны относительно отражения $A\mapsto-A$ и не
выбирают один из двух финалистов.
\end{theorem}

\begin{nogocriterion}[Запрет ориентации по знаку determinant]
Нельзя использовать знак determinant как физический селектор: он меняется
при переориентации базиса центральной плоскости. Нельзя также использовать
желаемую equal-gap мишень или меньший condition number как источник
конкретного $u/d$-ребра без parent-терма, различающего эти два типа.
\end{nogocriterion}

\begin{protocol}\begin{enumerate}
\item исходных rank-two назначений: \textbf{шесть};
\item maximum determinant survivors: \textbf{четыре};
\item maximum metric-area survivors: \textbf{четыре};
\item minimum-condition survivors: \textbf{два};
\item dimension-match survivors: \textbf{два};
\item финалисты: \textbf{$(u,Y)$ и $(d,Y)$};
\item знак determinant физически инвариантен: \textbf{нет};
\item gauge-неэквивалентность выбирает финалиста: \textbf{нет};
\item equal-gap мишень является внутренним селектором: \textbf{нет};
\item exact classification: \textbf{$8/8$};
\item intrinsic selector: \textbf{$0/8$};
\item следующий гейт: \textbf{parent-origin бинарного $u/d$-селектора}.
\end{enumerate}\end{protocol}

Машинный сертификат сохранён в
\path{s2t_v8_baryon_c0_singlet_triplet_central_gap_isotypic_channel_extra_edge_mass_two_source_cubic_incidence_assignment_selector_gate_results.json}.

\begin{thebibliography}{9}
\bibitem{v8-cubic-incidence-selector-krajewski} T.~Krajewski,
\emph{Classification of Finite Spectral Triples}, arXiv:hep-th/9701081.
\bibitem{v8-cubic-incidence-selector-cacic} B.~\'Caci\'c,
\emph{Moduli Spaces of Dirac Operators for Finite Spectral Triples},
arXiv:0902.2068.
\end{thebibliography}